% ============================================================
\section{Taller práctico: VAR para economía abierta}

% ============================================================
\begin{frame}{Taller práctico: economía abierta}

\begin{block}{Objetivo}
Construir escenarios macro-financieros para una economía pequeña y abierta utilizando un modelo VAR que incorpore variables domésticas y externas.
\end{block}

\begin{itemize}
    \item Evaluar la transmisión de shocks externos hacia la economía local.
    \item Simular trayectorias para alimentar ejercicios de stress testing.
    \item Comparar escenarios baseline, adverso externo, conjunto extremo y condicional.
\end{itemize}

\end{frame}

% ============================================================
\begin{frame}{Variables del modelo VAR}

\begin{columns}[T]

\begin{column}{0.48\textwidth}
\begin{block}{Variables domésticas}
\small
\begin{itemize}
    \item Crecimiento económico
    \item Inflación
    \item Tasa interbancaria
    \item Depreciación cambiaria
    \item EMBI
\end{itemize}
\end{block}
\end{column}

\begin{column}{0.48\textwidth}
\begin{block}{Variables externas}
\small
\begin{itemize}
    \item Precio internacional del petróleo
    \item Tasa de interés de la FED
    \item Crecimiento económico de EE.UU.
\end{itemize}
\end{block}
\end{column}

\end{columns}

\vspace{0.5em}

\begin{alertblock}{Idea central}
El modelo permite estudiar cómo un deterioro del entorno internacional se transmite a inflación, tipo de cambio, tasas, crecimiento y riesgo país.
\end{alertblock}

\end{frame}

% ============================================================
\begin{frame}{Narrativa del escenario adverso externo}

\begin{block}{Escenario}
A partir del inicio del horizonte de proyección, la economía enfrenta un deterioro del entorno internacional.
\end{block}

\begin{itemize}
    \item Aumento del precio internacional del petróleo.
    \item Incremento de la tasa de interés de la FED.
    \item Desaceleración del crecimiento de EE.UU.
    \item Presiones sobre el tipo de cambio.
    \item Aumento del EMBI.
    \item Mayor inflación y respuesta de la tasa doméstica.
    \item Menor crecimiento económico.
\end{itemize}

\end{frame}

% ============================================================
\begin{frame}{Calibración del shock adverso externo}

\begin{block}{Unidades}
Los shocks están expresados como múltiplos de la desviación estándar histórica de cada variable.
\end{block}

\begin{table}
\centering
\small
\begin{tabular}{lcc}
\toprule
Variable & Shock & Interpretación \\
\midrule
Crecimiento & $-1.5$ & menor actividad \\
Inflación & $+1.5$ & mayor presión inflacionaria \\
Interbancaria & $+1.0$ & respuesta monetaria \\
Depreciación FX & $+2.0$ & presión cambiaria \\
EMBI & $+2.0$ & mayor riesgo país \\
Precio petróleo & $+2.5$ & shock externo de costos \\
Tasa FED & $+2.0$ & condiciones financieras externas \\
Crecimiento EE.UU. & $-2.0$ & menor demanda externa \\
\bottomrule
\end{tabular}
\end{table}

\end{frame}

% ============================================================
\begin{frame}[fragile]{Estructura del script}

\small

\begin{verbatim}
# 1. Limpieza de sesión
# 2. Carga de paquetes
# 3. Parámetros generales
# 4. Carga y preparación de datos
# 5. Estimación del VAR
# 6. Escenario adverso externo
# 7. Escenario conjunto extremo
# 8. Escenario condicional
# 9. Tablas y gráficos
\end{verbatim}

\begin{block}{Producto final}
Un script reproducible en R que genere escenarios macro-financieros de economía abierta.
\end{block}

\end{frame}

% ============================================================
\begin{frame}[fragile]{Código base: variables del VAR}

\small

\begin{verbatim}
variables_var <- c(
  "crecimiento",
  "inflacion",
  "interbancaria",
  "deprec_fx",
  "embi",
  "oil_price",
  "fed_rate",
  "us_growth"
)
\end{verbatim}

\begin{alertblock}{Nota}
Las tres últimas variables representan el bloque externo del modelo.
\end{alertblock}

\end{frame}

% ============================================================
\begin{frame}[fragile]{Código base: estimación del VAR}

\small

\begin{verbatim}
Y <- macro %>%
  select(all_of(variables_var)) %>%
  na.omit()

modelo_var <- VAR(
  y = Y,
  p = 2,
  type = "const"
)

roots(modelo_var)
\end{verbatim}

\begin{block}{Diagnóstico}
Las raíces del VAR deben estar dentro del círculo unitario para que el modelo sea estable.
\end{block}

\end{frame}

% ============================================================
\begin{frame}[fragile]{Código base: shock adverso externo}

\small

\begin{verbatim}
shock_externo <- c(
  crecimiento   = -1.5,
  inflacion     =  1.5,
  interbancaria =  1.0,
  deprec_fx     =  2.0,
  embi          =  2.0,
  oil_price     =  2.5,
  fed_rate      =  2.0,
  us_growth     = -2.0
)

adverso_externo <- escenario_var_shock(
  modelo_var = modelo_var,
  h = h,
  shock = shock_externo
)
\end{verbatim}

\end{frame}

% ============================================================
\begin{frame}{Índice de severidad conjunta}

\begin{block}{Idea}
El escenario conjunto extremo no se obtiene tomando percentiles variable por variable.
\end{block}

\begin{enumerate}
    \item Se simulan muchas trayectorias conjuntas del VAR.
    \item Se calcula un índice de severidad para cada trayectoria.
    \item Se selecciona la trayectoria ubicada en la cola adversa.
\end{enumerate}

\vspace{0.5em}

\begin{alertblock}{Variables que aumentan la severidad}
Mayor inflación, tasa interbancaria, depreciación, EMBI, precio del petróleo y tasa FED; menor crecimiento local y menor crecimiento de EE.UU.
\end{alertblock}

\end{frame}

% ============================================================
\begin{frame}[fragile]{Código base: escenario condicional}

\small

\begin{verbatim}
restricciones <- tibble(
  crecimiento = c(1.0, 1.2, 1.5, 1.8, 2.0, rep(2.5, h - 5)),
  inflacion = c(8.0, 7.5, 7.0, 6.5, 6.0, rep(5.5, h - 5)),
  interbancaria = c(9.0, 9.0, 8.75, 8.5, 8.25, rep(8.0, h - 5)),
  deprec_fx = c(8.0, 7.0, 6.0, 5.0, 4.5, rep(4.0, h - 5)),
  embi = c(650, 625, 600, 575, 550, rep(525, h - 5)),
  oil_price = c(95, 92, 90, 88, 85, rep(82, h - 5)),
  fed_rate = c(6.0, 5.8, 5.5, 5.25, 5.0, rep(4.75, h - 5)),
  us_growth = c(0.5, 0.8, 1.0, 1.2, 1.5, rep(1.8, h - 5))
)
\end{verbatim}

\end{frame}

% ============================================================
\begin{frame}{Entregable del ejercicio}

\begin{block}{Producto final}
Una tabla consolidada con cinco períodos históricos y las trayectorias proyectadas para cada escenario.
\end{block}

\begin{table}
\centering
\small
\begin{tabular}{llll}
\toprule
Periodo & Escenario & Variable & Valor \\
\midrule
2025-08 & Histórico & Inflación & 4.2 \\
2025-09 & Histórico & Inflación & 4.5 \\
2026-01 & Baseline & Inflación & 4.8 \\
2026-01 & Adverso externo & Inflación & 7.5 \\
2026-01 & Condicional & Inflación & 8.0 \\
\bottomrule
\end{tabular}
\end{table}

\end{frame}

% ============================================================
\begin{frame}{Preguntas guía para los grupos}

\begin{enumerate}
    \item ¿Cuál es el shock externo dominante en el escenario?
    \item ¿Qué variable doméstica reacciona con mayor intensidad?
    \item ¿Cómo se transmite el shock del petróleo a la inflación?
    \item ¿Qué rol juega la tasa FED en el EMBI y el tipo de cambio?
    \item ¿Qué escenario luce más severo para el portafolio previsional?
\end{enumerate}

\end{frame}